\chapter{Le métier du trading de crypto-actifs}\label{ch:m3:the-crypto-trading-business}

Un projet de \omterm{def:m3:blockchains-for-traders:key}{jeton} sur le point d'être coté en bourse a besoin de quelqu'un pour le coter. Il prête à un teneur de marché deux pour cent
de son offre de \omterm{def:m3:blockchains-for-traders:key}{jetons} pour un an et lui accorde des calls portant sur le même montant, et le teneur de marché accepte
de maintenir un cours acheteur et un cours vendeur sur plusieurs plateformes. Aucune des parties ne verse de liquidités~; chacune pense avoir fait une bonne affaire. En
octobre 2024, le régulateur boursier américain a poursuivi trois sociétés se présentant comme des teneurs de
marché pour avoir vendu à des promoteurs de \omterm{def:m3:blockchains-for-traders:key}{jetons} une activité qui n'était pas du trading du tout~: des volumes échangés entre leurs
propres comptes, pour donner l'apparence de liquidité aux \omterm{def:m3:blockchains-for-traders:key}{jetons}. Ce chapitre porte sur l'activité qui entoure les marchés des
chapitres précédents~: les contrats en vertu desquels des sociétés font un marché de nouveaux \omterm{def:m3:blockchains-for-traders:key}{jetons} et la manière de les valoriser,
les desks de gré à gré qui négocient en grande taille, les opérations de conservation et de trésorerie qui protègent les crypto-actifs,
le fonctionnement d'un marché qui ne ferme jamais, et la réglementation qui arrive désormais.

\section{Contrats de tenue de marché de jetons}

\begin{definition}[Contrat de tenue de marché de jeton, structure prêt plus call]\label{def:m3:the-crypto-trading-business:tmma}
Un \emph{contrat de tenue de marché de jeton}\index{contrat de tenue de marché de jeton} est un contrat par lequel l'émetteur d'un \omterm{def:m3:blockchains-for-traders:key}{jeton} engage
une société pour fournir des cotations sur le \omterm{def:m3:blockchains-for-traders:key}{jeton}, sur des plateformes déterminées, avec des obligations portant sur le spread, la profondeur et la présence. Dans une
\emph{structure prêt plus call}\index{structure prêt plus call} l'émetteur paie le service en prêtant à la société des \omterm{def:m3:blockchains-for-traders:key}{jetons}
pour la durée du contrat, à restituer à la fin, et en lui accordant des options d'achat européennes sur des \omterm{def:m3:blockchains-for-traders:key}{jetons} à des prix d'exercice convenus.
\end{definition}

Le prêt donne au teneur de marché l'inventaire dont il a besoin pour offrir des \omterm{def:m3:blockchains-for-traders:key}{jetons} qu'il ne possède pas~; les calls constituent sa rémunération. Le
prêt seul n'entraîne aucune exposition nette pour le teneur de marché~: il détient les \omterm{def:m3:blockchains-for-traders:key}{jetons} et doit rendre le même nombre. Les calls
le rendent long du \omterm{def:m3:blockchains-for-traders:key}{jeton}. Un teneur de marché qui veut encaisser la valeur des calls plutôt que parier sur le \omterm{def:m3:blockchains-for-traders:key}{jeton} vend, dès le premier jour,
autant de \omterm{def:m3:blockchains-for-traders:key}{jetons} que le delta des calls (livre~1, chapitre~25), et ajuste cette couverture à mesure que le cours évolue~: il
vend lorsque le cours monte et rachète lorsqu'il baisse. La couverture est vendue sur le marché qui vient de s'ouvrir, à partir des
\omterm{def:m3:blockchains-for-traders:key}{jetons} prêtés.

\begin{proposition}[Valeur et couverture d'une opération prêt plus call]\label{prop:m3:the-crypto-trading-business:deal}
Soit l'émetteur qui prête $N$ \omterm{def:m3:blockchains-for-traders:key}{jetons} au cours $S$ pendant $T$ années et accorde des calls portant sur $n_i$ \omterm{def:m3:blockchains-for-traders:key}{jetons} aux prix d'exercice $K_i$. Avec une volatilité
$\sigma$ et sans dérive ni actualisation, le package vaut $\sum_i n_i C(S, K_i, T, \sigma)$, où $C$ est le prix du call de Black,
et sa valeur divisée par $NS$ est la commission implicite de l'opération. Le teneur de marché est neutre en delta après avoir vendu
$\sum_i n_i \Phi(d_{1,i})$ \omterm{def:m3:blockchains-for-traders:key}{jetons}, avec $d_{1,i} = \ln(S/K_i)/(\sigma\sqrt T) + \tfrac12\sigma\sqrt T$.
\end{proposition}

\begin{proof}
Les \omterm{def:m3:blockchains-for-traders:key}{jetons} reçus et les \omterm{def:m3:blockchains-for-traders:key}{jetons} dus se compensent, de sorte que la seule exposition est celle des calls. Leur valeur et leur delta, sous une loi lognormale sans dérive du cours,
sont donnés par la formule de Black et $\Phi(d_1)$ par \omterm{def:m3:blockchains-for-traders:key}{jeton} (livre~4 pour la version Monte Carlo utilisée dans le tutoriel).
\end{proof}

\begin{omfigure}
\centering
\begin{tikzpicture}[font=\small,
  p/.style={draw,thick,rounded corners=2pt,align=center,minimum height=1cm,text width=2.8cm},
  arr/.style={-{Stealth[length=2.2mm]},thick}]
\node[p,draw=omDef,fill=omDef!8,text width=2.2cm] (iss) at (0,0) {émetteur du jeton};
\node[p,draw=omThm,fill=omThm!8,text width=2.2cm] (mm) at (5.4,0) {teneur de marché};
\node[p,draw=omProp,fill=omProp!8,text width=2.2cm] (vx) at (10.2,0) {plateformes};
\draw[arr] ([yshift=8pt]iss.east) -- node[above,font=\footnotesize,align=center] {prêt de $N$ jetons,\\options d'achat} ([yshift=8pt]mm.west);
\draw[arr] ([yshift=-8pt]mm.west) -- node[below,font=\footnotesize,align=center] {cotations~; $N$ jetons\\restitués à la fin} ([yshift=-8pt]iss.east);
\draw[arr,{Stealth[length=2.2mm]}-{Stealth[length=2.2mm]}] (mm) -- node[above,font=\footnotesize] {cours acheteurs, cours vendeurs} node[below,font=\footnotesize] {couverture en delta} (vx);
\end{tikzpicture}

\omcaption{Une opération prêt plus call. L'émetteur paie avec un prêt de \omterm{def:m3:blockchains-for-traders:key}{jetons} et des options d'achat~; le teneur de marché cote sur les plateformes et couvre le
delta des calls en vendant une partie des \omterm{def:m3:blockchains-for-traders:key}{jetons} prêtés. Schéma.}\label{fig:m3:the-crypto-trading-business:deal}
\end{omfigure}

La structure porte un conflit intégré. Un teneur de marché long de calls gagne si le cours dépasse les prix d'exercice, et celui qui est payé
en \omterm{def:m3:blockchains-for-traders:key}{jetons} plutôt qu'en liquidités peut être tenté de l'y pousser~; un émetteur qui veut un prix de cotation élevé et un bon classement en volume peut souhaiter la même chose.
Les obligations mesurées sur le spread et la profondeur peuvent être tenues honnêtement~; celles mesurées sur le volume incitent au \omterm{def:m3:spot-markets:wash}{wash trading}
décrit au \cref{ch:m3:spot-markets}.

\begin{dated}{2026-09}{Manipulation de marché en tant que service}\label{dat:m3:the-crypto-trading-business:mirrors}
Le 9 octobre 2024, la SEC a annoncé des poursuites pour fraude contre trois sociétés se présentant comme des teneurs de marché (ZM Quant, Gotbit et CLS Global) et neuf
personnes physiques, au motif que des promoteurs de \omterm{def:m3:blockchains-for-traders:key}{jetons} avaient engagé ZM Quant et Gotbit pour fournir une «~manipulation de marché en tant que service~», incluant des volumes de transactions artificiels,
pour des \omterm{def:m3:blockchains-for-traders:key}{jetons} proposés à des investisseurs particuliers, et que ZM Quant et CLS Global avaient fait de même pour un \omterm{def:m3:blockchains-for-traders:key}{jeton} créé à la demande du FBI dans le cadre
d'une enquête parallèle. Le même jour, le bureau du procureur fédéral de Boston a levé le secret sur des poursuites pénales visant dix-huit personnes et entités,
dont quatre teneurs de marché, a saisi plus de 25~millions de dollars en crypto-actifs et désactivé des robots de trading qui avaient pratiqué du \omterm{def:m3:spot-markets:wash}{wash trading} sur environ 60 \omterm{def:m3:blockchains-for-traders:key}{jetons}.
\end{dated}

\begin{definition}[Déblocage de jetons]\label{def:m3:the-crypto-trading-business:unlock}
Un \emph{déblocage de jetons}\index{déblocage de jetons} est la libération programmée, à l'issue d'une période de blocage, de \omterm{def:m3:blockchains-for-traders:key}{jetons} alloués à l'équipe d'un projet, à ses investisseurs ou à sa trésorerie, après
quoi leurs détenteurs peuvent les vendre.
\end{definition}

Les déblocages ajoutent de l'offre à dates connues, et la couverture du teneur de marché ainsi que ses cotations doivent en tenir compte~: un déblocage important est un vendeur futur connu,
comme une sortie d'indice l'est dans le livre~1, chapitre~15.

\begin{omfigure}
\begin{tikzpicture}
\begin{axis}[width=11cm,height=5cm,
  xlabel={volatilité du jeton, \% par an},ylabel={commission implicite, \% du prêt},
  tick label style={font=\small},label style={font=\small},
  xmin=40,xmax=240,ymin=0,ymax=70,grid=major,grid style={gray!25},
  table/col sep=comma]
\addplot[omThm,very thick] table[x=vol,y=fee]{figdata/markets-3/24-the-crypto-trading-business/fee.csv};
\end{axis}
\end{tikzpicture}

\omcaption{La commission implicite de l'opération du tutoriel (2~\% de l'offre prêtée pour un an à un cours de lancement de 0.50~\$, calls portant chacun sur un tiers aux prix d'exercice de 0.75,
1.00 et 1.50) en fonction de la volatilité du \omterm{def:m3:blockchains-for-traders:key}{jeton} (\cref{prop:m3:the-crypto-trading-business:deal}). À 120~\% la valeur des calls représente 26~\% du prêt.
Opération illustrative. Données~: le tutoriel du chapitre.}\label{fig:m3:the-crypto-trading-business:fee}
\end{omfigure}

\begin{omfigure}
\begin{tikzpicture}
\begin{axis}[width=11cm,height=5cm,
  xlabel={cours du jeton, \$},ylabel={couverture, millions de jetons vendus à découvert},
  tick label style={font=\small},label style={font=\small},
  xmin=0.1,xmax=3,ymin=0,ymax=2.1,grid=major,grid style={gray!25},
  legend columns=2,legend style={font=\footnotesize,at={(0.5,-0.3)},anchor=north,draw=none,fill=none,/tikz/every even column/.append style={column sep=8pt}},
  table/col sep=comma]
\addplot[omThm,very thick] table[x=price,y=t1]{figdata/markets-3/24-the-crypto-trading-business/hedge.csv};\addlegendentry{à la signature (un an restant)}
\addplot[omDef,very thick,dashed] table[x=price,y=t05]{figdata/markets-3/24-the-crypto-trading-business/hedge.csv};\addlegendentry{six mois restants}
\end{axis}
\end{tikzpicture}

\omcaption{Couverture en delta du teneur de marché sur l'opération de cotation~: \omterm{def:m3:blockchains-for-traders:key}{jetons} vendus à découvert face aux calls, en fonction du cours du \omterm{def:m3:blockchains-for-traders:key}{jeton}. Il vend lorsque le cours
monte (vers les 2~millions de \omterm{def:m3:blockchains-for-traders:key}{jetons} des calls) et rachète lorsqu'il baisse. Données~: le tutoriel du chapitre.}\label{fig:m3:the-crypto-trading-business:hedge}
\end{omfigure}

\section{Desks OTC}

Les transactions de grande taille en crypto-actifs se font hors bourse, par demande de cotation (RFQ) (livre~2, chapitre~22)~: un fonds demande à un desk un prix pour 500 bitcoins, et le desk
cote et, si elle est exécutée, se couvre sur les bourses en quelques minutes ou quelques heures. Le prix du desk est le coût espéré de la couverture (la profondeur qu'elle consommera sur l'ensemble des plateformes, les frais,
les jambes en \omterm{def:m3:blockchains-for-traders:stable}{stablecoins} et en monnaie fiduciaire), plus le risque que le cours varie pendant la couverture, plus une marge. Le règlement-livraison est la partie délicate~: la monnaie fiduciaire transite par les banques
aux heures bancaires, tandis que les crypto-actifs circulent sur les chaînes à toute heure, si bien qu'un côté de la transaction est toujours en transit. Les desks règlent contre des soldes préfinancés, ou livrent en premier
aux clients disposant de lignes de crédit, et le risque de crédit de la fenêtre de règlement est tarifé comme le risque de règlement des marchés de change (livre~2, chapitre~20).

\section{Trésorerie et opérations de conservation}

\begin{definition}[Auto-conservation, portefeuille chaud, portefeuille froid, calcul multipartite]\label{def:m3:the-crypto-trading-business:custody}
L'\emph{auto-conservation}\index{auto-conservation} consiste à détenir des crypto-actifs sous des \omterm{def:m3:blockchains-for-traders:key}{clés privées} que le propriétaire contrôle, plutôt que comme solde auprès d'une plateforme ou d'un dépositaire. Un
\emph{portefeuille chaud}\index{portefeuille chaud} est un portefeuille dont les clés se trouvent sur des systèmes connectés à internet, capables de signer immédiatement~; un \emph{portefeuille froid}\index{portefeuille froid} est celui
dont les clés sont conservées hors ligne, de sorte que la signature suppose un processus délibéré hors ligne. Le \emph{calcul multipartite}\index{calcul multipartite} (MPC) est une
méthode cryptographique par laquelle plusieurs parties, détenant chacune une part d'une clé, produisent conjointement une signature sans que la clé n'existe jamais en un seul lieu.
\end{definition}

La trésorerie d'une société de trading conserve l'essentiel de ses crypto-actifs à froid et ne transfère vers les \omterm{def:m3:the-crypto-trading-business:custody}{portefeuilles chauds} et les plateformes que ce dont l'activité de trading a besoin, avec une approbation par MPC ou multisignature afin
qu'aucune personne seule ni aucun serveur ne puisse déplacer les fonds. Les défaillances sont anciennes, sous des habits neufs. FTX conservait presque tous ses crypto-actifs dans des \omterm{def:m3:the-crypto-trading-business:custody}{portefeuilles chauds}, sans contrôles multisignatures
(\cref{ch:m3:centralised-exchanges}). Chez QuadrigaCX, une bourse canadienne effondrée en 2019, le personnel de l'Ontario Securities Commission a constaté que le
cofondateur et directeur général s'était crédité de soldes fictifs sous des pseudonymes, les avait négociés contre des clients, et avait perdu davantage sur des
plateformes extérieures~; au total, plus de 76\,000 clients ont perdu au moins 169 millions de dollars canadiens. Le personnel a qualifié l'affaire d'«~une fraude à l'ancienne, enveloppée dans une technologie
moderne~». La séparation des tâches, le rapprochement quotidien des soldes on-chain avec les livres et les limites sur les personnes habilitées à approuver un transfert constituent les garde-fous, comme
dans toute autre classe d'actifs.

\begin{omfigure}
\centering
\begin{tikzpicture}[font=\small,
  w/.style={draw,thick,rounded corners=2pt,align=center,minimum height=1cm,text width=2.5cm},
  arr/.style={-{Stealth[length=2.2mm]},thick}]
\node[w,draw=omDef,fill=omDef!8] (c) at (0,0) {conservation à froid\\(l'essentiel des actifs)};
\node[w,draw=omThm,fill=omThm!8] (m) at (3.6,0) {MPC ou multisignature\\approbations};
\node[w,draw=omExo,fill=omExo!8] (h) at (7.2,0) {portefeuilles chauds\\(besoins quotidiens)};
\node[w,draw=omProp,fill=omProp!8] (v) at (10.8,0) {plateformes, OTC\\contreparties};
\draw[arr] (c) -- (m); \draw[arr] (m) -- (h); \draw[arr] (h) -- (v);
\draw[arr,dashed] (v.south) to[bend left=20] node[below,font=\footnotesize] {rapatriement des soldes excédentaires} (c.south);
\end{tikzpicture}

\omcaption{La chaîne de conservation d'une société de trading~: les actifs ne sortent de la conservation à froid que par une approbation multipartite, les \omterm{def:m3:the-crypto-trading-business:custody}{portefeuilles chauds} détenant ce dont la journée a besoin, et les soldes excédentaires
sur les plateformes sont rapatriés. Schéma de principe.}\label{fig:m3:the-crypto-trading-business:custody}
\end{omfigure}

\section{Opérations 24 heures sur 24}

Les marchés de crypto-actifs ne ferment jamais. Une société a besoin de systèmes et de personnes pour chaque heure~: surveillance des positions, des marges et du collatéral sur chaque plateforme, couverture des
week-ends où les banques sont fermées mais où les crypto-actifs circulent, procédures pour les mises à niveau des chaînes et les fenêtres de maintenance des plateformes, et des ingénieurs d'astreinte pour la connectivité décrite au
\cref{ch:m3:crypto-data-and-infrastructure}. La liquidité d'un week-end est plus mince et ses mouvements plus amples rapportés à la profondeur~; les liquidations
traitées au \cref{ch:m3:margin-liquidation-and-loss-allocation} se produisent à n'importe quel moment. Les limites de risque de la société doivent être appliquées par des machines qui ne dorment jamais, avec des opérateurs
alertés par exception, et sa liquidité en monnaie fiduciaire doit être positionnée avant la fermeture des banques le vendredi.

\section{Régulation}

\begin{definition}[Prestataire de services sur crypto-actifs, règlement sur les marchés de crypto-actifs]\label{def:m3:the-crypto-trading-business:mica}
Un \emph{prestataire de services sur crypto-actifs}\index{prestataire de services sur crypto-actifs} (PSCA, crypto-asset service provider) est une société qui fournit des services sur crypto-actifs tels que la conservation, l'exploitation d'une
plateforme de négociation, l'échange de crypto-actifs, l'exécution ou la réception et la transmission d'ordres, pour lesquels le droit de l'Union européenne exige un agrément. Le \emph{règlement sur les marchés de
crypto-actifs}\index{règlement sur les marchés de crypto-actifs} (Markets in Crypto-Assets Regulation, MiCA) est le règlement de l'Union européenne, règlement (UE) 2023/1114, qui fixe des règles pour
les émetteurs de crypto-actifs et de \omterm{def:m3:blockchains-for-traders:stable}{stablecoins} et pour les PSCA.
\end{definition}

\begin{dated}{2026-09}{Calendrier de MiCA}\label{dat:m3:the-crypto-trading-business:mica}
L'ESMA indique que MiCA est entré en vigueur en juin 2023~; ses titres sur les \omterm{def:m3:blockchains-for-traders:key}{jetons} se référant à des actifs et sur les \omterm{def:m3:blockchains-for-traders:key}{jetons} de monnaie électronique (\omterm{def:m3:blockchains-for-traders:stable}{stablecoins}) s'appliquent à compter du 30 juin 2024, et le
reste à compter du 30 décembre 2024~; en vertu de son article 143(3), les entités fournissant déjà des services sur crypto-actifs peuvent poursuivre leur activité jusqu'au 1\up{er} juillet 2026 ou jusqu'à ce qu'elles soient agréées ou
refusées.
\end{dated}

\begin{definition}[Travel rule du GAFI]\label{def:m3:the-crypto-trading-business:travel}
La \emph{travel rule}\index{travel rule} est l'exigence, figurant dans la recommandation 16 du GAFI sur les virements, selon laquelle les informations sur le donneur d'ordre et le
bénéficiaire accompagnent un virement entre entreprises réglementées~; le GAFI a étendu ses normes aux actifs virtuels et à leurs prestataires de services en 2018.
\end{definition}

La note interprétative du GAFI de juin 2019 applique la recommandation 16 aux transferts d'actifs virtuels et permet aux pays d'exempter les transferts en dessous d'un seuil
ne dépassant pas 1\,000~USD/EUR. L'Union européenne n'en a retenu aucun~: son règlement 2023/1113, applicable à partir du 30 décembre 2024, exige que soient transmis
avec chaque transfert de crypto-actifs, quel qu'en soit le montant, le nom du donneur d'ordre, son adresse sur le registre et ses éléments d'identification.

\begin{dated}{2026-09}{La loi américaine sur les stablecoins}\label{dat:m3:the-crypto-trading-business:genius}
Le GENIUS Act, Public Law 119-27 du 18 juillet 2025, organise la régulation des \omterm{def:m3:blockchains-for-traders:stable}{stablecoins} de paiement~; il impose à tout émetteur autorisé de maintenir des
réserves identifiables sur une base d'au moins un pour un, en pièces et en billets, en dépôts à vue, en bons du Trésor (T-bill) d'une échéance de 93 jours ou moins, en pensions livrées et
en actifs similaires.
\end{dated}

Pour une société de trading, la régulation modifie la carte du \cref{ch:m3:centralised-exchanges}~: quelles plateformes elle peut utiliser et depuis où, si ses activités propres (conservation
pour le compte de clients, tenue de marché pour les émetteurs, négociation de gré à gré) exigent un agrément, ce qu'elle doit déclarer et à qui, et quelles informations doivent
accompagner ses transferts.

\section{Tutoriel~: valorisation d'une opération de cotation}\label{tut:m3:the-crypto-trading-business:deal}

\begin{tutorial}
\textbf{Objectif.} Valoriser par simulation un prêt de \omterm{def:m3:blockchains-for-traders:key}{jetons} assorti de trois strikes de calls, en déduire les frais qu'il implique, et calculer la couverture que le teneur de marché doit détenir au lancement
et tout au long de l'année. \textbf{État final~:} les \cref{fig:m3:the-crypto-trading-business:fee,fig:m3:the-crypto-trading-business:hedge} et les chiffres du
problème du week-end.
\begin{tutsteps}
\item \textbf{Le package.} Monte-Carlo avec paires antithétiques et le cours terminal comme variable de contrôle, vérifié par rapport à Black.
\omcode{code/firm/tokenloan/firm_tokenloan.py}{39}{55}{Valeur Monte-Carlo du package de calls avec variable de contrôle.}\label{lst:m3:the-crypto-trading-business:mc}
\item \textbf{Frais et couverture.} \texttt{implied\_fee}, \texttt{day\_one\_hedge} et \texttt{delta\_schedule}.
\omcode{code/firm/tokenloan/firm_tokenloan.py}{62}{74}{Les frais implicites, la couverture du premier jour et la couverture tout au long de l'année.}\label{lst:m3:the-crypto-trading-business:hedge}
\item \textbf{Lancer} \texttt{summary()} et \texttt{fig\_business.py}.
\end{tutsteps}
\textbf{Pour aller plus loin.} Ajouter un \omterm{def:m3:the-crypto-trading-business:unlock}{déblocage de jetons} de 10~\% de l'offre au sixième mois, sous forme de saut d'offre qui fait bouger le cours~; imputer à la couverture la fourchette qu'elle paie
sur un marché peu liquide~; valoriser l'opération du côté de l'émetteur, y compris les \omterm{def:m3:blockchains-for-traders:key}{jetons} qu'il récupère.
\end{tutorial}

\section{Construction~: contrats de tenue de marché de jetons}\label{bld:m3:the-crypto-trading-business:tokenloan}

\begin{build}
\bfield{Rôle} La société miniature signe, valorise et couvre des \omterm{def:m3:the-crypto-trading-business:tmma}{contrats de tenue de marché de jetons}, surveille ses obligations de cotation et comptabilise la valeur qu'elle gagne.
\bfield{Interface} \texttt{Deal(tokens\_loaned, price, years, vol, tranches)}~; \texttt{call\_bs}~; \texttt{package\_value\_mc}, \texttt{package\_value}~;
\texttt{implied\_fee}~; \texttt{day\_one\_hedge}~; \texttt{delta\_schedule(deal, prices, t\_left)}~; \texttt{obligations\_met(spread\_bp, depth\_usd, uptime)}~;
\texttt{post\_fee(ledger, deal, ts)}~: écritures dans \texttt{firm.ledger}.
\bfield{Règles} \omterm{def:m3:blockchains-for-traders:key}{Jetons} reçus et dus compensés~; volatilité tirée de l'historique propre du \omterm{def:m3:blockchains-for-traders:key}{jeton} et des cotations comparables, jamais supposée faible~; obligations vérifiées par plateforme
et par jour~; les calls comptabilisés à la signature en produits de frais dans le grand livre.
\bfield{Tests d'acceptation} \texttt{code/firm/tokenloan/tests/}~: Monte-Carlo à 1~\% près de la forme fermée~; limites des calls~; frais compris entre 0 et 1~; une couverture qui augmente avec le
cours~; obligations~; une écriture dans le grand livre.
\bfield{Pour aller plus loin} Calendriers de déblocage~; coûts de couverture sur des carnets d'ordres peu profonds~; comptabilisation du passif du prêt~; valeur de l'opération pour l'émetteur.
\end{build}

\begin{omsources}
\item SEC, communiqué de presse 2024-166, 9 octobre 2024.
\item ESMA, Markets in Crypto-Assets Regulation (MiCA) page; Regulation (EU) 2023/1114. Consulté en septembre 2026.
\item GENIUS Act, Public Law 119-27, 18 juillet 2025 (GovInfo).
\item Commission des valeurs mobilières de l'Ontario, \emph{QuadrigaCX: A Review by Staff of the Ontario Securities Commission}, 2020.
\item GAFI, \emph{International Standards on Combating Money Laundering and the Financing of Terrorism and Proliferation} (les Recommandations du GAFI),
mise à jour de juin 2025~: Recommandations 15 et 16 et Note interprétative de la Recommandation 15 (copie sur Internet Archive).
\item Regulation (EU) 2023/1113 of 31 May 2023 on information accompanying transfers of funds and certain crypto-assets (copie d'EUR-Lex sur Internet Archive).
\item US Attorney's Office, District of Massachusetts, communiqué de presse du 9 octobre 2024 (copie sur Internet Archive).
\end{omsources}

\section{Exercices}

\begin{exercise}[$\star$]\label{exo:m3:the-crypto-trading-business:1}
Expliquer pourquoi un prêt de \omterm{def:m3:blockchains-for-traders:key}{jetons}, à lui seul, ne laisse au teneur de marché aucune exposition nette au \omterm{def:m3:blockchains-for-traders:key}{jeton}.
\end{exercise}

\begin{exercise}[$\star$]\label{exo:m3:the-crypto-trading-business:2}
Distinguer un \omterm{def:m3:the-crypto-trading-business:custody}{portefeuille chaud} d'un \omterm{def:m3:the-crypto-trading-business:custody}{portefeuille froid}, et dire ce qu'apporte le \omterm{def:m3:the-crypto-trading-business:custody}{calcul multipartite} (MPC).
\end{exercise}

\begin{exercise}[$\star$]\label{exo:m3:the-crypto-trading-business:3}
Indiquer à partir de quand MiCA s'applique aux émetteurs de \omterm{def:m3:blockchains-for-traders:stable}{stablecoins}, et à partir de quand aux \omterm{def:m3:the-crypto-trading-business:mica}{prestataires de services sur crypto-actifs}.
\end{exercise}

\begin{exercise}[$\star\star$]\label{exo:m3:the-crypto-trading-business:4}
Évaluer un call sur un \omterm{def:m3:blockchains-for-traders:key}{jeton} à 0.50~USD, strike 0.75, un an, volatilité 120~\%, sans dérive ni actualisation.
\end{exercise}

\begin{exercise}[$\star\star$]\label{exo:m3:the-crypto-trading-business:5}
Expliquer comment les frais implicites de l'opération d'admission à la cote varient entre des volatilités de 80~\% et de 160~\%, et pourquoi cette variation est si forte.
\end{exercise}

\begin{exercise}[$\star\star$]\label{exo:m3:the-crypto-trading-business:6}
Pour un client OTC qui demande un prix sur 500 bitcoins, énumérer les éléments qui entrent dans la cotation du desk.
\end{exercise}

\begin{exercise}[$\star\star\star$]\label{exo:m3:the-crypto-trading-business:7}
\emph{Programmation.} Avec \texttt{delta\_schedule}, calculer le nombre de \omterm{def:m3:blockchains-for-traders:key}{jetons} que le teneur de marché doit racheter si le cours passe de 0.50 à 0.25 avec six mois restants.
\end{exercise}

\begin{exercise}[$\star\star\star$]\label{exo:m3:the-crypto-trading-business:8}
\emph{Repérer l'erreur.} «~Notre accord ne nous rémunère que si le volume sur le \omterm{def:m3:blockchains-for-traders:key}{jeton} atteint 5~millions de dollars par jour~; nous veillerons à ce qu'il le fasse.~»
\end{exercise}

\section{Problème~: L'opération d'admission à la cote}

\begin{problem}[{Problème du week-end --- ce que valent les calls et ce qu'il faut vendre}]\label{pb:m3:the-crypto-trading-business:1}
Un projet de 100 millions de \omterm{def:m3:blockchains-for-traders:key}{jetons} est admis à la cote à 0.50~USD. Il prête à un teneur de marché 2~\% de l'offre pendant un an et accorde des calls portant chacun sur un tiers de ce montant, à des prix d'exercice
de 0.75, 1.00 et 1.50~USD. La volatilité du \omterm{def:m3:blockchains-for-traders:key}{jeton} est estimée à 120~\% par an~; négliger les intérêts et la dérive.

\medskip\noindent\textbf{Partie~I --- La valeur.}
\begin{enumerate}
\item Calculer la valeur du prêt en dollars.
\item Calculer la valeur des trois tranches de calls, puis celle du package.
\item Déterminer les frais que le package implique, exprimés en proportion du prêt.
\item Mesurer l'écart entre la valeur Monte-Carlo et la forme fermée.
\item Calculer les frais pour une volatilité de 80~\%.
\end{enumerate}

\medskip\noindent\textbf{Partie~II --- La couverture.}
\begin{enumerate}[resume]
\item Calculer le nombre de \omterm{def:m3:blockchains-for-traders:key}{jetons} que le teneur de marché vend dès le premier jour.
\item Déterminer la part du prêt que cela représente.
\item Décrire l'effet sur la couverture si le \omterm{def:m3:blockchains-for-traders:key}{jeton} double au cours du premier mois.
\item Faire de même s'il est divisé par deux.
\item Expliquer pourquoi la vente du premier jour compte pour le nouveau marché.
\end{enumerate}

\medskip\noindent\textbf{Partie~III --- Les obligations.}
\begin{enumerate}[resume]
\item Indiquer quelles obligations l'accord doit fixer, et lesquelles il doit éviter.
\item Expliquer comment le teneur de marché respecterait une fourchette maximale de 1~\% avec une profondeur de 50\,000~USD de chaque côté.
\item Indiquer l'effet d'un \omterm{def:m3:the-crypto-trading-business:unlock}{déblocage de jetons} au sixième mois sur la couverture et sur les cotations.
\item Expliquer comment la société doit comptabiliser l'opération.
\item Identifier ce qui rendrait cette opération illégale.
\end{enumerate}

\medskip\noindent\textbf{Partie~IV --- Jugement critique.}
\begin{enumerate}[resume]
\item Dire quelle partie tire le meilleur parti de l'opération.
\item Expliquer pourquoi les émetteurs paient en calls plutôt qu'en numéraire.
\item Indiquer comment s'assurer qu'un teneur de marché est honnête.
\item Énoncer le \emph{résultat nommé}~: la valeur du package de calls en proportion de la valeur des \omterm{def:m3:blockchains-for-traders:key}{jetons} prêtés, et le nombre de \omterm{def:m3:blockchains-for-traders:key}{jetons} vendus le premier jour.
\item En une phrase, dire ce que vend à l'émetteur un teneur de marché lors d'une admission à la cote.
\end{enumerate}
\end{problem}

\section{Questions d'entretien}

\begin{interviewq}[$\star$ \iqroles{trader}]\label{iq:m3:the-crypto-trading-business:1}
Comment un teneur de marché est-il rémunéré dans une \omterm{def:m3:the-crypto-trading-business:tmma}{structure prêt plus call}, et comment se couvre-t-il~?
\end{interviewq}

\begin{interviewq}[$\star$ \iqroles{risque}]\label{iq:m3:the-crypto-trading-business:2}
Décrivez un dispositif de conservation solide pour une société de trading sur crypto-actifs.
\end{interviewq}

\begin{interviewq}[$\star\star$ \iqroles{chercheur}]\label{iq:m3:the-crypto-trading-business:3}
Comment estimeriez-vous la volatilité d'un \omterm{def:m3:blockchains-for-traders:key}{jeton} qui n'a pas encore été négocié~?
\end{interviewq}

\begin{interviewq}[$\star\star$ \iqroles{trader}]\label{iq:m3:the-crypto-trading-business:4}
Comment évaluez-vous un bloc de gré à gré de 500 bitcoins à 3~h du matin un dimanche~?
\end{interviewq}

\begin{interviewq}[$\star\star$ \iqroles{risque}]\label{iq:m3:the-crypto-trading-business:5}
Qu'est-ce qui change pour une société de trading lorsque MiCA et la \omterm{def:m3:the-crypto-trading-business:travel}{travel rule} s'y appliquent~?
\end{interviewq}

\begin{interviewq}[$\star\star\star$ \iqroles{développeur, risque}]\label{iq:m3:the-crypto-trading-business:6}
Concevez les contrôles pour transférer 50~millions de dollars de crypto-actifs depuis les \omterm{def:m3:the-crypto-trading-business:custody}{portefeuilles froids} vers trois plateformes de négociation en moins d'une heure.
\end{interviewq}
